% Source: Weinberg GR, physical PDF pp. 172--173
%         (printed pp. 147--148).
% Coverage: Section 6.9, The Geometric Analogy; equation (6.9.1).
% Figures/tables/footnotes: optional-reading footnote; no figures or tables.
% Status: source-reviewed and compile-clean.
% Uncertainties: none.

\subsection{The Geometric Analogy\optionalreading}\label{sec:6.9}

\begin{quote}
\small\emph{This section lies somewhat out of the book's main line of
development and may be omitted in a first reading.}
\end{quote}

We have seen in this chapter that the nonvanishing of the tensor
\(R_{\lambda\mu\nu\kappa}\) is the true expression of the presence of a
gravitational field. We also saw in Chapter~1 that Gauss was led to introduce
the Gaussian curvature \(K=R/2\) as the true measure of the departure of a
two-dimensional geometry from that of Euclid, and that Riemann subsequently
introduced the curvature tensor \(R_{\lambda\mu\nu\kappa}\) to generalize the
concept of curvature to three or more dimensions. It is therefore not
surprising that Einstein and his successors have regarded the effects of a
gravitational field as producing a change in the geometry of space and time.
At one time it was even hoped that the rest of physics could be brought into
a geometric formulation, but this hope has met with disappointment, and the
geometric interpretation of the theory of gravitation has dwindled to a mere
analogy, which lingers in our language in terms like ``metric,'' ``affine
connection,'' and ``curvature,'' but is not otherwise very useful. The
important thing is to be able to make predictions about images on the
astronomers' photographic plates, frequencies of spectral lines, and so on,
and it simply does not matter whether we ascribe these predictions to the
physical effect of gravitational fields on the motion of planets and photons
or to a curvature of space and time. (The reader should be warned that these
views are heterodox and would meet with objections from many general
relativists.)

Despite the preceding remarks, it is worth mentioning without proof just what
the tensor \(R_{\lambda\mu\nu\kappa}\) has to do with the curvature of a
Riemannian space. Given a point \(X\) in a space of an arbitrary number of
dimensions, and given two vectors \(a^\mu,b^\mu\) defined at \(X\), we may
construct through \(X\) a family of ``geodesic curves''
\(x^\mu=x^\mu(\tau,\alpha,\beta)\) defined by
\[
  \frac{\dd^2x^\mu}{\dd\tau^2}
  +\Gamma^\mu{}_{\nu\lambda}
   \frac{\dd x^\nu}{\dd\tau}
   \frac{\dd x^\lambda}{\dd\tau}
  =0,
  \qquad
  \left.\frac{\dd x^\mu}{\dd\tau}\right|_{x=X}
  =
  \alpha a^\mu+\beta b^\mu,
\]
the numbers \(\alpha,\beta\) being allowed to run over all real values. These
curves fill out a two-dimensional surface \(S(a,b)\) through \(X\), and the
Gaussian curvature of this surface at \(X\)
is\hyperref[ch6-ref-5]{\textsuperscript{5}}
\begin{equation}
  K(a,b)
  =
  \frac{
    R_{\lambda\mu\nu\kappa}
    a^\lambda b^\mu a^\nu b^\kappa
  }{
    (g_{\lambda\nu}g_{\mu\kappa}
     -g_{\lambda\kappa}g_{\mu\nu})
    a^\lambda b^\mu a^\nu b^\kappa
  }.
  \tag{6.9.1}\label{eq:6.9.1}
\end{equation}
From Eq.~\eqref{eq:6.7.4} we see that in two dimensions \(K(a,b)\) is
independent of \(a\) and \(b\), and is just \(R/2\).
